% !TEX root = ../main.tex
\begin{ouverture}[spectre]{2}{Absorbance et dosage par étalonnage}{Partie 1 · Constitution et transformations de la matière}
{Couleur\enspace·\enspace Spectre d'absorption\enspace·\enspace Absorbance\enspace·\enspace Loi de Beer-Lambert\enspace·\enspace Gamme étalon}
\objectif{Expliquer ou prévoir la couleur d'une espèce en solution à partir de son spectre UV-visible.}
\objectif{Exploiter la loi de Beer-Lambert et connaître ses conditions de validité.}
\objectif{Déterminer la concentration d'un soluté à partir de l'absorbance de solutions de concentrations connues.}
\objectif{Proposer un protocole pour réaliser une gamme étalon et tester les limites du dosage.}
\prerequis{Lumière blanche, spectres, longueur d'onde (classe de seconde).}
\prerequis{Dosage par étalonnage, dilution et concentrations (classe de seconde et chapitre~1).}
\prerequis{Lecture graphique et proportionnalité.}
\end{ouverture}
\essentiel
\essSection{1}{La couleur d'une espèce en solution}
\begin{definition}[Couleur perçue]
Éclairée en lumière blanche, une solution absorbe certaines radiations et transmet les autres. Lorsqu'elle absorbe
surtout un domaine de longueurs d'onde, sa couleur perçue est la \textbf{couleur complémentaire} de la couleur absorbée :
sur le cercle chromatique, elles sont diamétralement opposées. Une solution qui n'absorbe aucune radiation visible est
incolore.
\end{definition}
\schema{02-cercle}
\essSection{2}{Absorbance et spectre d'absorption}
\begin{definition}[Absorbance]
Un spectrophotomètre envoie sur la solution une lumière monochromatique d'intensité $I_0$ et mesure l'intensité $I$ qu'elle
transmet. L'\textbf{absorbance} $A=\log\dfrac{I_0}{I}$, sans unité, est d'autant plus grande que la solution absorbe. Le
\textbf{spectre d'absorption} représente $A$ en fonction de la longueur d'onde $\lambda$ ; on repère la longueur d'onde $\lambda_{\max}$ du
maximum d'absorption.
\end{definition}
\schema{02-spectrophotometre}
\schema{02-spectre}
\essSection{3}{La loi de Beer-Lambert}
\begin{loi}[Loi de Beer-Lambert]
Pour une solution diluée d'une espèce absorbante, à une longueur d'onde donnée, l'absorbance est proportionnelle à la
concentration de cette espèce :
\[A=\varepsilon\,\ell\,c=k\,c,\]
où $\ell$ est l'épaisseur de solution traversée (en cm), $c$ la concentration (en \unit{mol.L^{-1}}) et $\varepsilon$ le \textbf{coefficient
d'absorption molaire} (en \unit{L.mol^{-1}.cm^{-1}}), qui dépend de l'espèce, de la longueur d'onde et du solvant. La loi est valable
pour une lumière monochromatique et des solutions assez diluées, en pratique tant que $A$ reste inférieure à 2 environ.
Si plusieurs espèces absorbent, leurs absorbances s'ajoutent.
\end{loi}
\begin{exemple}
Pour le permanganate de potassium à \qty{525}{nm}, $\varepsilon=\qty{2,4e3}{L.mol^{-1}.cm^{-1}}$. Dans une cuve de \qty{1,00}{cm}, une solution à \qty{1,0e-4}{mol.L^{-1}}
a pour absorbance $A=\num{2.4e3}\times1{,}00\times\num{1.0e-4}=0{,}24$.
\end{exemple}
\begin{attention}
L'absorbance n'a pas d'unité. Dans $A=\varepsilon\,\ell\,c$, les unités doivent être cohérentes : $\ell$ en cm si $\varepsilon$ est en \unit{L.mol^{-1}.cm^{-1}}. On
ne lit une concentration que \textbf{dans} le domaine de la gamme étalon : au-delà, rien ne garantit que la loi s'applique.
\end{attention}
\essSection{4}{Le dosage par étalonnage}
\begin{savoirfaire}[Doser une espèce colorée par étalonnage]
\begin{questions}
\item Tracer le spectre de l'espèce et régler le spectrophotomètre sur $\lambda_{\max}$, où la mesure est la plus sensible ;
faire le « blanc » avec le solvant.
\item Préparer une \textbf{gamme étalon} : des solutions de concentrations connues, obtenues par dilution d'une solution mère.
\item Mesurer leurs absorbances et tracer $A$ en fonction de $c$ : les points s'alignent sur une droite passant par l'origine,
ce qui confirme la loi de Beer-Lambert.
\item Mesurer l'absorbance de la solution à doser, diluée si nécessaire pour rester dans le domaine de la gamme, puis lire
ou calculer sa concentration, sans oublier le facteur de dilution.
\end{questions}
\end{savoirfaire}
\schema{02-etalonnage}
\begin{codecote}[0.62]
\begin{pythoncol}
c = [1.0, 2.0, 4.0, 6.0, 8.0]        # µmol/L
A = [0.106, 0.212, 0.424, 0.636, 0.848]
num, den = 0, 0
for i in range(len(c)):
    num = num + c[i] * A[i]
    den = den + c[i] ** 2
k = num / den
print(round(k, 3), round(0.45 / k, 2))
\end{pythoncol}
\begin{consolecol}
0.106 4.25
\end{consolecol}
\cote
\fontsize{9.5}{12.5}\selectfont
\capacite{Capacité numérique}\par\smallskip
Le programme calcule le coefficient directeur $k$ de la droite passant par l'origine la plus proche des points de la
gamme, par la méthode des moindres carrés : $k=\dfrac{\sum c_i\,A_i}{\sum c_i^2}$. Il en déduit la concentration de la solution
S, d'absorbance 0,45 : \qty{4,25}{\micro mol.L^{-1}}.
\end{codecote}
\begin{experience}[Doser le bleu patenté d'un sirop de menthe]
À partir d'une solution mère de bleu patenté à \qty{10,0}{\micro mol.L^{-1}}, on prépare une gamme de cinq solutions filles. Après
avoir fait le blanc avec de l'eau distillée, on mesure leurs absorbances à \qty{638}{nm}, puis celle du sirop dilué vingt fois.
La droite d'étalonnage donne la concentration du sirop dilué, que l'on multiplie par 20.
\end{experience}
\begin{rediger}
Justifier le choix de la longueur d'onde ($\lambda_{\max}$) et la validité de la loi (points alignés sur une droite passant par
l'origine). Montrer la lecture graphique par des pointillés ou calculer $c=\dfrac{A}{k}$. Conclure en tenant compte du
facteur de dilution, avec l'unité demandée.
\end{rediger}
\exercices

\rubrique{Découverte}
\begin{blocLivre}{ex-2-1}
\begin{exercice}[1]{Spectre et couleur}
\schema{02-ex01}
\begin{questions}
\item Déterminer la longueur d'onde du maximum d'absorption de chaque solution.
\item Quelle couleur chaque solution absorbe-t-elle le plus ?
\item À l'aide du cercle chromatique, prévoir la couleur perçue de chaque solution.
\item La solution 2 est une solution de permanganate de potassium. Quelle est sa couleur ?
\end{questions}
\end{exercice}
\end{blocLivre}
\begin{blocLivre}{ex-2-2}
\begin{exercice}[1]{Prévoir l'allure d'un spectre}
\begin{questions}
\item Une solution de sulfate de cuivre paraît bleu-cyan. Dans quel domaine de longueurs d'onde absorbe-t-elle ?
\item Une solution de diiode paraît jaune-orangé. Quelle couleur absorbe-t-elle ?
\item Une solution de chlorophylle paraît verte ; son spectre présente deux maximums. Où se situent-ils approximativement ?
\item Une solution n'absorbe aucune radiation visible. Quelle est sa couleur ? Peut-on la doser par spectrophotométrie dans
le visible ?
\end{questions}
\end{exercice}
\end{blocLivre}
\begin{blocLivre}{ex-2-3}
\begin{exercice}[1]{Appliquer la loi de Beer-Lambert}
Pour le permanganate de potassium, $\ce{KMnO4}$, à \qty{525}{nm} : $\varepsilon=\qty{2,4e3}{L.mol^{-1}.cm^{-1}}$. La cuve a une épaisseur $\ell=\qty{1,00}{cm}$.
\begin{donnees}
Masses molaires atomiques (\unit{g.mol^{-1}}) : O 16,0 ; K 39,1 ; Mn 54,9.
\end{donnees}
\begin{questions}
\item Calculer l'absorbance d'une solution de concentration \qty{2,0e-4}{mol.L^{-1}}.
\item Une autre solution a pour absorbance 0,30. Calculer sa concentration.
\item Que devient l'absorbance de la première solution dans une cuve de \qty{2,00}{cm} ?
\item Calculer la concentration en masse de la solution de la question 2, en \unit{mg.L^{-1}}.
\end{questions}
\end{exercice}
\end{blocLivre}
\begin{blocLivre}{ex-2-4}
\begin{exercice}[1]{Lire une droite d'étalonnage}
On a mesuré à \qty{800}{nm} l'absorbance de solutions de sulfate de cuivre de concentrations connues.
\schema{02-ex04}
\begin{questions}
\item Les mesures vérifient-elles la loi de Beer-Lambert ? Justifier.
\item Déterminer le coefficient $k$ de la relation $A=k\,c$, puis le coefficient d'absorption molaire $\varepsilon$.
\item Une solution S de sulfate de cuivre a pour absorbance 0,54. Déterminer sa concentration.
\item En déduire la concentration en masse de sulfate de cuivre de S, avec $M(\ce{CuSO4})=\qty{159,6}{g.mol^{-1}}$.
\end{questions}
\end{exercice}
\end{blocLivre}
\begin{blocLivre}{ex-2-5}
\begin{exercice}[1]{Diluer avant de mesurer}
Un sirop de menthe pur a une absorbance de 3,1 à \qty{638}{nm}. Dilué dix fois, il a une absorbance de 0,42. Pour le bleu
patenté à \qty{638}{nm}, dans la même cuve, $A=k\,c$ avec $k=\qty{0,106}{L.\micro mol^{-1}}$.
\begin{questions}
\item Pourquoi la mesure sur le sirop pur n'est-elle pas exploitable ?
\item Calculer la concentration en bleu patenté du sirop dilué, puis celle du sirop.
\item Décrire la préparation de \qty{100,0}{mL} de sirop dilué dix fois.
\end{questions}
\end{exercice}
\end{blocLivre}
\rubrique{Entraînement}
\begin{blocLivre}{ex-2-6}
\begin{exercice}[2]{Préparer une gamme étalon}
On dispose d'une solution mère S$_0$ de bleu patenté à $c_0=\qty{10,0}{\micro mol.L^{-1}}$. On veut préparer, dans des fioles jaugées de
\qty{25,0}{mL}, cinq solutions filles de concentrations 1,0 ; 2,0 ; 4,0 ; 6,0 et \qty{8,0}{\micro mol.L^{-1}}.
\begin{questions}
\item Établir l'expression du volume $V_0$ de solution mère à prélever pour obtenir une solution fille de concentration $c$.
\item Calculer les cinq volumes à prélever.
\item Quel instrument permet de prélever tous ces volumes avec précision ?
\item Pourquoi faut-il mesurer toutes les absorbances à la même longueur d'onde et avec la même cuve ?
\end{questions}
\end{exercice}
\end{blocLivre}
\begin{blocLivre}{ex-2-7}
\begin{exercice}[2]{Le permanganate de l'eau de Dakin}
L'eau de Dakin, un antiseptique, contient de l'hypochlorite de sodium, incolore, et du permanganate de potassium, qui lui
donne sa couleur rose. Son étiquette indique « permanganate de potassium : 0,0010~\% », soit \qty{0,0010}{g} pour \qty{100}{mL}.
On mesure à \qty{525}{nm}, dans une cuve de \qty{1,00}{cm}, une absorbance de 0,150.
\begin{donnees}
$\ce{KMnO4}$ : $M=\qty{158,0}{g.mol^{-1}}$ ; $\varepsilon=\qty{2,4e3}{L.mol^{-1}.cm^{-1}}$ à \qty{525}{nm}.
\end{donnees}
\begin{questions}
\item Pourquoi l'eau de Dakin est-elle rose ? Pourquoi mesure-t-on à \qty{525}{nm} ?
\item L'hypochlorite de sodium gêne-t-il la mesure ? Justifier.
\item Calculer la concentration en permanganate, puis sa concentration en masse.
\item Convertir l'indication de l'étiquette en \unit{mg.L^{-1}} et conclure.
\end{questions}
\end{exercice}
\end{blocLivre}
\begin{blocLivre}{ex-2-8}
\begin{exercice}[2]{Le cuivre d'une bouillie bordelaise}
La bouillie bordelaise, utilisée pour traiter la vigne, contient du sulfate de cuivre. L'étiquette annonce « 20~\% de cuivre »
en masse. On dissout \qty{1,00}{g} de produit dans une fiole jaugée de \qty{100,0}{mL} ; tout le cuivre passe en solution sous forme
d'ions $\ce{Cu^{2+}}$. L'absorbance de la solution à \qty{800}{nm} vaut 0,37.
\begin{donnees}
Étalonnage : $A=k\,c$ avec $k=\qty{0,012}{L.mmol^{-1}}$ (exercice 4) ; $M(\ce{Cu})=\qty{63,5}{g.mol^{-1}}$.
\end{donnees}
\begin{questions}
\item Calculer la concentration en ions cuivre de la solution.
\item En déduire la masse de cuivre contenue dans l'échantillon.
\item Calculer le pourcentage massique de cuivre du produit. Le résultat est-il conforme à l'étiquette ?
\item Pourquoi faut-il vérifier que l'absorbance mesurée est comprise dans le domaine de la gamme étalon ?
\end{questions}
\end{exercice}
\end{blocLivre}
\begin{blocLivre}{ex-2-9}
\begin{exercice}[2]{Tester la loi de Beer-Lambert en Python}
Un élève a mesuré l'absorbance de cinq solutions de bleu patenté et teste la proportionnalité avec ce programme.
\begin{python}
c = [1.0, 2.0, 4.0, 6.0, 8.0]           # µmol/L
A = [0.107, 0.210, 0.425, 0.634, 0.851]
num = 0
den = 0
for i in range(len(c)):
    num = num + c[i] * A[i]
    den = den + c[i] ** 2
k = num / den
print(round(k, 4))
for i in range(len(c)):
    print(c[i], round(100 * (A[i] - k * c[i]) / (k * c[i]), 1))
\end{python}
\begin{console}
0.1061
1.0 0.8
2.0 -1.0
4.0 0.1
6.0 -0.4
8.0 0.3
\end{console}
\begin{questions}
\item \capacite{Capacité numérique}Que représente la valeur 0.1061 ? Quelle est son unité ?
\item Que calcule la seconde boucle ? Interpréter les valeurs affichées.
\item La loi de Beer-Lambert est-elle vérifiée ? Justifier.
\item Une solution inconnue a pour absorbance 0,50. Quelle ligne ajouter pour afficher sa concentration ? Quelle valeur
obtient-on ?
\end{questions}
\end{exercice}
\end{blocLivre}
\begin{blocLivre}{ex-2-10}
\begin{exercice}[2]{Les limites de la loi de Beer-Lambert}
On mesure l'absorbance de solutions de bleu patenté de plus en plus concentrées.
\schema{02-ex10}
\begin{questions}
\item Sur quel domaine de concentrations la loi de Beer-Lambert est-elle vérifiée ? Justifier.
\item Quelle absorbance prévoirait la loi pour \qty{25}{\micro mol.L^{-1}} ? Comparer à la mesure.
\item Une solution inconnue a une absorbance de 2,4. Peut-on en déduire sa concentration avec la droite ? Proposer une
dilution adaptée.
\end{questions}
\end{exercice}
\end{blocLivre}
\begin{blocLivre}{ex-2-11}
\begin{exercice}[2]{Un sirop vert, deux colorants}
Un sirop de menthe contient deux colorants : la tartrazine (E102), jaune, et le bleu patenté (E131). On a tracé le spectre
du sirop dilué vingt fois.
\schema{02-ex11}
\begin{questions}
\item Expliquer la couleur verte du sirop à partir de son spectre.
\item Quel colorant est responsable de chacun des deux maximums ?
\item À \qty{638}{nm}, la tartrazine n'absorbe pas. Pourquoi peut-on doser le bleu patenté à cette longueur d'onde ?
\item Calculer la concentration en bleu patenté du sirop dilué, avec $k=\qty{0,106}{L.\micro mol^{-1}}$ à \qty{638}{nm}.
\end{questions}
\end{exercice}
\end{blocLivre}
\rubrique{Objectif DS}
\begin{blocLivre}{ex-2-12}
\begin{exercice}[3]{Combien de verres de sirop ?}
\begin{docu}{1}{Le colorant bleu}
Le bleu patenté V (E131) est un colorant alimentaire autorisé. Sa dose journalière admissible (DJA) est de \qty{5}{mg} par
kilogramme de masse corporelle et par jour (Autorité européenne de sécurité des aliments, 2013).
\end{docu}
\begin{docu}{2}{Le dosage}
Un sirop de menthe est dilué vingt fois. L'absorbance de la solution diluée, mesurée à \qty{638}{nm} dans une cuve de \qty{1,00}{cm},
vaut 0,53. La droite d'étalonnage du bleu patenté a pour équation $A=0{,}106\,c$, avec $c$ en \unit{\micro mol.L^{-1}}. Un verre de
boisson contient \qty{20}{mL} de sirop.
\end{docu}
\begin{donnees}
Masse molaire du bleu patenté (sel de sodium) : \qty{582,7}{g.mol^{-1}}.
\end{donnees}
\begin{questions}
\item Calculer la concentration en bleu patenté de la solution diluée, puis celle du sirop.
\item En déduire la concentration en masse du sirop en bleu patenté, en \unit{mg.L^{-1}}.
\item Calculer la masse de bleu patenté contenue dans un verre.
\item \initiative Un enfant de \qty{30}{kg} risque-t-il de dépasser la DJA en buvant du sirop de menthe ?
\end{questions}
\end{exercice}
\end{blocLivre}
\begin{blocLivre}{ex-2-13}
\begin{exercice}[3]{Les nitrates d'une eau de puits}
\begin{docu}{1}{Une limite de qualité}
Pour être déclarée potable, une eau doit contenir moins de \qty{50}{mg.L^{-1}} d'ions nitrate, $\ce{NO3-}$ (Code de la santé publique).
\end{docu}
\begin{docu}{2}{Le dosage}
Les ions nitrate, incolores, sont transformés par un réactif en une espèce jaune dont on mesure l'absorbance à \qty{410}{nm}.
\par\smallskip\noindent\hfill\begin{tabular}{lcccccc}
\toprule
$t(\ce{NO3-})$ (\unit{mg.L^{-1}}) & 0 & 10 & 20 & 30 & 40 & 50 \\
$A$ & 0 & 0,085 & 0,170 & 0,255 & 0,340 & 0,425 \\
\bottomrule
\end{tabular}\hfill\null\par\smallskip
Traitées de la même manière, l'eau du puits A donne une absorbance de 0,392, celle du puits B une absorbance de 0,48.
\end{docu}
\begin{donnees}
N 14,0 ; O 16,0 (\unit{g.mol^{-1}}).
\end{donnees}
\begin{questions}
\item Pourquoi faut-il transformer les ions nitrate avant de mesurer une absorbance dans le visible ?
\item Les mesures de la gamme vérifient-elles la loi de Beer-Lambert ? Déterminer le coefficient de proportionnalité.
\item Déterminer la concentration en masse des ions nitrate de l'eau du puits A, puis sa concentration en quantité de matière.
\item \initiative Les eaux des deux puits sont-elles potables vis-à-vis des nitrates ? Discuter la fiabilité de la conclusion
pour le puits B.
\end{questions}
\end{exercice}
\end{blocLivre}
\begin{blocLivre}{ex-2-14}
\begin{exercice}[3]{La caféine d'une boisson énergisante}
\begin{docu}{1}{L'étiquette}
Une boisson énergisante indique : « caféine : \qty{32}{mg} pour \qty{100}{mL} ». Une canette contient \qty{250}{mL}.
\end{docu}
\begin{docu}{2}{Le dosage dans l'ultraviolet}
La caféine, incolore, absorbe dans l'ultraviolet ; son spectre présente un maximum à \qty{273}{nm}, où $\varepsilon=\qty{9,7e3}{L.mol^{-1}.cm^{-1}}$.
La boisson, dégazée, est diluée cinquante fois ; l'absorbance de la solution diluée à \qty{273}{nm}, dans une cuve de \qty{1,00}{cm} en
quartz, vaut 0,318.
\end{docu}
\begin{docu}{3}{Une recommandation}
Pour un adolescent, l'Autorité européenne de sécurité des aliments considère qu'un apport de caféine de \qty{3}{mg} par
kilogramme de masse corporelle et par jour est sans risque (2015).
\end{docu}
\begin{donnees}
Caféine : $\ce{C8H10N4O2}$, $M=\qty{194,2}{g.mol^{-1}}$.
\end{donnees}
\begin{questions}
\item Calculer l'absorbance qu'aurait la boisson non diluée. Pourquoi faut-il la diluer ?
\item Calculer la concentration en caféine de la solution diluée, puis celle de la boisson.
\item En déduire la concentration en masse de la boisson et comparer à l'étiquette.
\item \initiative Combien de canettes un adolescent de \qty{60}{kg} peut-il boire par jour sans dépasser la recommandation ?
\end{questions}
\end{exercice}
\end{blocLivre}
\begin{blocLivre}{ex-2-15}
\begin{exercice}[3]{Le fer d'un complément alimentaire}
\begin{docu}{1}{Le comprimé}
Un complément alimentaire annonce \qty{14}{mg} de fer par comprimé, sous forme d'ions $\ce{Fe^{2+}}$.
\end{docu}
\begin{docu}{2}{Le protocole}
On dissout un comprimé dans une fiole jaugée de \qty{1,000}{L}. On prélève \qty{2,00}{mL} de cette solution, que l'on introduit dans une
fiole jaugée de \qty{50,0}{mL} avec un réactif : les ions $\ce{Fe^{2+}}$ forment une espèce rouge-orangé. Après ajustement au trait de
jauge, l'absorbance à \qty{510}{nm}, dans une cuve de \qty{1,00}{cm}, vaut 0,108. À cette longueur d'onde, $\varepsilon=\qty{1,11e4}{L.mol^{-1}.cm^{-1}}$ ; chaque ion
$\ce{Fe^{2+}}$ donne une entité colorée.
\end{docu}
\begin{donnees}
$M(\ce{Fe})=\qty{55,8}{g.mol^{-1}}$.
\end{donnees}
\begin{questions}
\item Quel est le facteur de dilution entre la fiole de \qty{1,000}{L} et celle de \qty{50,0}{mL} ?
\item Calculer la concentration de l'espèce colorée dans la fiole de \qty{50,0}{mL}.
\item \initiative Déterminer la masse de fer contenue dans le comprimé et la comparer à l'indication de l'étiquette.
\item Quelle absorbance aurait-on mesurée si le comprimé contenait exactement \qty{14}{mg} de fer ?
\end{questions}
\end{exercice}
\end{blocLivre}
